\begin{apendicesenv}

\chapter{MAPEAMENTO DO RELATO AO PRISMA 2020}
\label{apendice:prisma-checklist}

Este apêndice apresenta um mapeamento do relato deste manuscrito com base no PRISMA 2020 \cite{PRISMA2020}. A tabela desdobra os 27 itens do checklist em suas alíneas aplicáveis e utiliza tópicos abreviados; os requisitos completos estão no \href{https://www.prisma-statement.org/prisma-2020-checklist}{checklist oficial}. O mapeamento avalia o que foi apresentado no texto, não certifica a execução do procedimento, o rigor metodológico, a ausência de viés ou a certeza da evidência. São registradas as informações localizadas e as limitações do manuscrito, sem atribuir pontuação ou declarar conformidade global ou integral.

A indicação \textit{Localizado} identifica uma passagem pertinente, mas não certifica sua suficiência. \textit{Parcial} indica informação disponível com limitações explicitadas. \textit{Não apresentado} registra informação ou resultado não encontrado no manuscrito examinado, sem afirmar que o procedimento jamais ocorreu; quando a atividade não foi realizada, isso é indicado na justificativa. \textit{Não aplicável} é empregado somente com justificativa específica; a ausência de meta-análise não dispensa o relato dos métodos de síntese narrativa.

\begingroup
\small
\setlength{\tabcolsep}{4pt}
\renewcommand{\arraystretch}{1.12}
\begin{longtable}{|>{\raggedright\arraybackslash}p{0.8cm}|>{\raggedright\arraybackslash}p{3.0cm}|>{\raggedright\arraybackslash}p{10.0cm}|}
\caption{Mapeamento do relato aos itens do PRISMA 2020.}
\label{tab:prisma-checklist}\\
\hline
\textbf{Item} & \textbf{Tópico abreviado} & \textbf{Localização e estado do relato}\\
\hline
\endfirsthead
\multicolumn{3}{c}{\tablename\ \thetable\ -- Continuação}\\
\hline
\textbf{Item} & \textbf{Tópico abreviado} & \textbf{Localização e estado do relato}\\
\hline
\endhead
\hline
\multicolumn{3}{r}{\textit{Continua na próxima página}}\\
\endfoot
\hline
\endlastfoot
1 & Título & Parcial: capa e folha de rosto identificam o tema, mas o título não explicita revisão sistemática. \\
\hline
2 & Resumo & Parcial: Resumo e Abstract apresentam método e resultados; não foi concluída a conferência específica do PRISMA para resumos. \\
\hline
3 & Justificativa & Localizado: Cap.~\ref{cap:introducao}, Justificativa. \\
\hline
4 & Objetivos & Localizado: Cap.~\ref{cap:introducao}, Problema de Pesquisa e Objetivos. \\
\hline
5 & Elegibilidade & Parcial: Cap.~\ref{cap:metodologia}, Critérios de Seleção. Os critérios estão descritos, mas os 18 retidos permanecem uma população operacional: 17 são candidatos empíricos provisórios e 1 é protocolo/proposta contextual, fora da síntese de evidências. \\
\hline
6 & Fontes & Parcial: Cap.~\ref{cap:metodologia}, Bases de Dados e APIs. Não estão discriminadas as datas da última consulta de cada fonte. \\
\hline
7 & Buscas & Parcial: Cap.~\ref{cap:metodologia}, Estratégia de Busca Bilíngue. Há composição canônica e exemplos; o relato por fonte, com filtros e limites, não está completo. \\
\hline
8 & Seleção & Parcial: Cap.~\ref{cap:metodologia}, Processo de Seleção, e Cap.~\ref{cap:revisao}, Limitações da Revisão. São descritos automação e revisor único; a verificação documental permanece incompleta. \\
\hline
9 & Coleta de dados & Parcial: Cap.~\ref{cap:metodologia}, Identificação e Avaliação da Qualidade Metodológica. A obtenção de metadados está descrita; a extração dos textos e sua confirmação não estão integralmente documentadas. \\
\hline
10a & Desfechos & Parcial: Cap.~\ref{cap:metodologia}, Critérios de Seleção, e Tabela~\ref{tab:sintese-estudos-empiricos}. Faltam regras completas para seleção de resultados por desfecho. \\
\hline
10b & Outras variáveis & Parcial: Cap.~\ref{cap:metodologia}, Identificação. Não há definição completa das variáveis extraídas e do tratamento de informação ausente. \\
\hline
11 & Viés nos estudos & Parcial: Cap.~\ref{cap:metodologia}, Avaliação da Qualidade Metodológica. O MMAT está documentado como apreciação preliminar por revisor único, com fontes e julgamentos ainda incompletos. \\
\hline
12 & Medidas de efeito & Parcial: Tabela~\ref{tab:sintese-estudos-empiricos}. As métricas reportadas são heterogêneas; não foi definida medida comum para combinação quantitativa. \\
\hline
13a & Composição das sínteses & Parcial: Cap.~\ref{cap:metodologia}, Elegibilidade, e Cap.~\ref{cap:revisao}, Síntese dos Estudos Incluídos. A separação entre os 17 candidatos empíricos provisórios e o protocolo contextual está relatada; o procedimento para formar cada síntese e a confirmação científica final dos candidatos não estão completamente detalhados. \\
\hline
13b & Preparação dos dados & Não apresentado: não foi localizado procedimento específico para preparar os resultados extraídos, incluindo estatísticas ausentes ou conversões. \\
\hline
13c & Apresentação & Localizado: Cap.~\ref{cap:revisao}, Fluxo de Seleção, Estatísticas Descritivas, Síntese dos Estudos Incluídos e Avaliação Metodológica com o MMAT. A apresentação não substitui a validação dos dados. \\
\hline
13d & Síntese & Parcial: Cap.~\ref{cap:revisao}, Síntese dos Estudos Incluídos e Limitações da Revisão. A heterogeneidade é discutida; a codificação e o procedimento da síntese narrativa não estão inteiramente explicitados. \\
\hline
13e & Heterogeneidade & Não apresentado: não foi realizado método formal para investigar possíveis causas de heterogeneidade; diferenças entre populações, desenhos e métricas são discutidas qualitativamente. \\
\hline
13f & Sensibilidade & Não apresentado: não foi realizada análise de sensibilidade da síntese. \\
\hline
14 & Viés de relato & Não apresentado: não foi realizada avaliação formal. Cap.~\ref{cap:revisao}, Limitações da Revisão, registra a ausência de teste de viés de publicação. \\
\hline
15 & Certeza & Não apresentado: não foi apresentada avaliação formal da certeza do conjunto de evidências; o MMAT individual não a substitui. \\
\hline
16a & Fluxo & Parcial: Figura~\ref{fig:prisma-flow} e Tabela~\ref{tab:stats-descritivas}. Há contagens operacionais, mas não a discriminação completa de relatórios procurados, não recuperados e examinados. \\
\hline
16b & Exclusões justificadas & Não apresentado: o manuscrito não traz uma relação de estudos aparentemente elegíveis excluídos, com respectivas citações e razões. O ledger de adjudicação versionado registra decisões de escopo, mas não substitui essa apresentação no TCC. \\
\hline
17 & Características & Parcial: Tabela~\ref{tab:sintese-estudos-empiricos}. Os candidatos são citados, mas características de amostra e desenho não estão uniformemente detalhadas. \\
\hline
18 & Apreciação individual & Parcial: Seção~\ref{sec:qualidade-estudos} e Tabela~\ref{tab:mmat-reavaliacao-atual}. Julgamentos preliminares, com limitações de fontes e localizadores. \\
\hline
19 & Resultados individuais & Parcial: Tabela~\ref{tab:sintese-estudos-empiricos}. Diversas entradas descrevem a abordagem; estimativas e precisão não são apresentadas uniformemente. \\
\hline
20a & Síntese e limitações & Parcial: Cap.~\ref{cap:revisao}, Síntese dos Estudos Incluídos e Avaliação Metodológica com o MMAT. Não está completa a ligação entre os estudos de cada síntese e seus julgamentos metodológicos. \\
\hline
20b & Sínteses estatísticas & Não aplicável aos resultados de combinação estatística: nenhuma síntese estatística conjunta foi realizada. \\
\hline
20c & Heterogeneidade & Não apresentado: não foram apresentados resultados de investigação formal das causas de heterogeneidade. \\
\hline
20d & Sensibilidade & Não apresentado: não foram apresentados resultados de análises de sensibilidade. \\
\hline
21 & Viés de relato & Não apresentado: não foram apresentados resultados de avaliação formal de viés decorrente de resultados ausentes. \\
\hline
22 & Certeza & Não apresentado: não foi apresentada classificação formal de certeza por desfecho. \\
\hline
23a & Interpretação & Parcial: Cap.~\ref{cap:resultados-discussao}, Discussão Integrada. As interpretações dependem da consolidação da extração e da apreciação metodológica. \\
\hline
23b & Limites da evidência & Localizado: Cap.~\ref{cap:revisao}, Avaliação Metodológica com o MMAT e Limitações da Revisão; Cap.~\ref{cap:conclusao}, Limitações. \\
\hline
23c & Limites do processo & Localizado: Cap.~\ref{cap:metodologia}, Avaliação da Qualidade Metodológica, e Cap.~\ref{cap:revisao}, Limitações da Revisão, incluindo revisor único e acesso desigual aos textos. \\
\hline
23d & Implicações & Localizado: Cap.~\ref{cap:resultados-discussao}, Discussão Integrada, e Cap.~\ref{cap:conclusao}, Continuidade da Pesquisa. \\
\hline
24a & Registro & Localizado: Cap.~\ref{cap:metodologia}, Protocolo e Relato, declara ausência de registro prospectivo. \\
\hline
24b & Protocolo & Parcial: Cap.~\ref{cap:metodologia} descreve o planejamento; não identifica endereço de acesso a uma versão completa do protocolo anterior à execução. \\
\hline
24c & Alterações & Não apresentado: não foi localizada uma declaração metodológica completa de alterações do protocolo, com justificativas. A ausência de registro prospectivo não responde, por si só, a este item. \\
\hline
25 & Apoio & Não apresentado: não foi localizada declaração específica sobre apoio financeiro ou não financeiro e papel dos apoiadores na revisão. \\
\hline
26 & Interesses & Não apresentado: não foi localizada declaração específica de conflitos de interesse; não se presume sua ausência. \\
\hline
27 & Materiais & Parcial: Cap.~\ref{cap:metodologia}, Reprodutibilidade, descreve código e exports. A relação de materiais públicos e respectivos endereços precisa ser explicitada. \\
\hline
\end{longtable}
\endgroup

A apreciação MMAT corresponde aos itens 11 e 18 deste mapeamento e permanece preliminar. Ela não constitui avaliação formal da certeza do conjunto de evidências (itens 15 e 22), não é avaliação de viés de relato (itens 14 e 21) e não permite inferir ausência de viés ou conformidade global. As pendências indicadas restringem o alcance das conclusões e devem ser consideradas juntamente com as limitações apresentadas no Capítulo~\ref{cap:revisao}.

\end{apendicesenv}
